\documentclass{article}
\usepackage{hyperref}
\usepackage{amsmath}
\usepackage{calc}           % for \widthof in dimension expressions
\usepackage{propositions}

\DeclareNumberedStyle{P}

% A wide-label style for propformlabel demos
\SetPropStyle{widelevelone}{
  counter     = numpropi,
  format      = Item~(#1),
}

\begin{document}

% ============================================================
\section{Env-level \texttt{style=}}
% ============================================================

% style= at env level sets \l__props_env_style_tl.  This is used in the
% style-resolution step for every \item that doesn't give its own style=.

\subsection{All items thesis style}

% The two lists below should look identical.

Via env-level \texttt{style=thesis}:
\begin{prop}[style=thesis]
  \item[Hypothesis A] First. \label{ea:A}
  \item[Hypothesis B] Second. \label{ea:B}
\end{prop}

Via per-item \texttt{style=thesis}:
\begin{prop}
  \item[style=thesis, Hypothesis A] First.
  \item[style=thesis, Hypothesis B] Second.
\end{prop}

Refs: \ref{ea:A}, \ref{ea:B}.

\subsection{Env style for nameless items}

% style= at env level is used whether or not a name is given.
% Here all nameless items use the widelevelone style instead of
% the usual levelone.

\begin{prop}[style=widelevelone]
  \item First (should be ``Item~(1)'').     \label{eb:1}
  \item Second (should be ``Item~(2)'').    \label{eb:2}
\end{prop}

Refs: \ref{eb:1}, \ref{eb:2}.

\subsection{Per-item \texttt{style=} overrides env style}

\begin{prop}[style=thesis]
  \item[Hypothesis C] Thesis style (from env).
  \item[style=proposition, Override] Proposition style (per item overrides env).
\end{prop}

% ============================================================
\section{Env-level item-key defaults}
% ============================================================

% Non-style keys in the env arg are collected into \l__props_env_item_keys_tl
% and applied in phase 2.5 (BEFORE style_load).  This means the style can
% override them.  Keys that the style does not set survive to display time.

\subsection{Important ordering note}

% The style resolution step only sees keys set in phase 1 (the item's own
% optional arg), not the env defaults (phase 2.5).  So:
%
%   \begin{prop}[name=foo]\item bar
%
% resolves to the NAMELESS style (because name is empty in phase 1), then
% applies name=foo in phase 2.5.  Result: name "foo" formatted as "(foo)"
% by the levelone style.  The counter is NOT stepped (name non-empty).
%
%   \begin{prop}\item[name=foo] bar
%
% resolves to the NAMED style (proposition, bold) because name is set in
% phase 1.  Result: \textbf{foo}.
%
% These are different!  To get the same result as item-level name=foo,
% combine with an explicit env-level style=:

``\verb|\begin{prop}[name=foo]|'' (nameless style chosen; counter not stepped):
\begin{prop}[name=foo]
  \item bar
\end{prop}

``\verb|\begin{prop}\item[name=foo]|'' (named style; proposition bold):
\begin{prop}
  \item[name=foo] bar
\end{prop}

``\verb|\begin{prop}[style=proposition, name=foo]|'' (same as item-level):
\begin{prop}[style=proposition, name=foo]
  \item bar
\end{prop}

\subsection{counter= as env default (combined with style=)}

% counter= as an env default survives style load when the style does not
% set counter.  Use style=proposition (which sets only format, not counter).

Via env-level \texttt{counter=P} and \texttt{style=proposition}:
\begin{prop}[style=proposition, counter=P]
  \item \label{ec:P1}  % steps P, displays bold P-number
  \item \label{ec:P2}
\end{prop}

Via per-item \texttt{counter=P}:
\begin{prop}
  \item[counter=P] \label{ec:P3}
  \item[counter=P] \label{ec:P4}
\end{prop}

Refs: \ref{ec:P1}, \ref{ec:P2}, \ref{ec:P3}, \ref{ec:P4}.

\subsection{Per-item key overrides env default}

\begin{prop}[style=proposition, counter=P]
  \item First P. \label{ec2:first}
  \item[counter=numpropi] This item uses numpropi instead. \label{ec2:np}
  \item Back to P. \label{ec2:back}
\end{prop}

Refs: \ref{ec2:first} (P), \ref{ec2:np} (numpropi), \ref{ec2:back} (P).

% ============================================================
\section{\texttt{propformlabel} and \texttt{items=}}
% ============================================================

% \propformlabel is computed at \begin{prop} (and \begin{inlineprop})
% from a preview item resolution with the current env settings.
% It is updated after each \item to reflect the most recently output label.
%
% The items= key says how many items to advance the counter before
% computing the preview label, so propformlabel reflects the label
% that will appear after approximately that many items.
%
% Because \propformlabel is set before the list dimensions are applied,
% it can be used directly in leftmargin, labelwidth, etc.

\subsection{Basic use: leftmargin from first label}

% With a numbered style, propformlabel at open is the label for item 1.
% leftmargin=\widthof{\propformlabel}+0.5em sizes the margin to that label.

Preview label is ``\propformlabel'' (before any prop).

\begin{prop}[leftmargin=\widthof{\propformlabel}+0.5em]
  \item At open, label previewed as ``\propformlabel''. \label{pfl:1}
  \item After first item, \texttt{\string\propformlabel} = ``\propformlabel''. \label{pfl:2}
  \item After second item, \texttt{\string\propformlabel} = ``\propformlabel''. \label{pfl:3}
\end{prop}

\subsection{items= for a longer list: size for last expected label}

% When items=20, the preview label is what would appear as the 20th item:
% e.g. "(20)".  This gives enough left margin for the widest label.

\begin{prop}[items=20, leftmargin=\widthof{\propformlabel}+0.5em]
  \item Margin is sized for ``\propformlabel'' (shown at open). \label{pfl:long1}
  \item Second item. \label{pfl:long2}
  \item Third item. \label{pfl:long3}
\end{prop}

After the list the counter has advanced by 3 (not 20).

\subsection{propformlabel with a named style}

% propformlabel also works with named styles that use counters.

\begin{prop}[style=proposition, counter=P, items=5,
             leftmargin=\widthof{\propformlabel}+0.5em]
  \item First P with sized margin.   \label{pfl:P1}
  \item Second P with sized margin.  \label{pfl:P2}
\end{prop}

Refs: \ref{pfl:P1}, \ref{pfl:P2}.

\subsection{propformlabel with widelevelone}

% Shows that propformlabel adapts to the style's format.
% With widelevelone, the label is ``Item (n)''; items=9 previews ``Item (9)''.

\begin{prop}[style=widelevelone, items=9,
             leftmargin=\widthof{\propformlabel}+0.5em]
  \item First (margin sized for ``Item~(9)''). \label{pfl:w1}
  \item Second. \label{pfl:w2}
\end{prop}

\subsection{Using propformlabel with labelindent}

% labelindent places the left edge of the label at a given distance
% from the enclosing margin.  Here we size it to zero and let leftmargin
% be determined by the label width.

\begin{prop}[style=widelevelone, items=5,
             labelindent=0em,
             leftmargin=\widthof{\propformlabel}+\labelsep]
  \item First (flush label, margin = label width + labelsep). \label{pfl:li1}
  \item Second. \label{pfl:li2}
\end{prop}

% ============================================================
\section{\texttt{\textbackslash propoptions} item-key defaults}
% ============================================================

% Item-level keys set via \propoptions persist across environments.
% They are group-scoped: \begingroup ... \endgroup scopes the effect.
% Env-arg keys override propoptions keys for that environment.

\subsection{Persistent counter= default}

% \propoptions{counter=P, style=proposition} applies to all following envs.
% Both envs below should produce bold P-numbered items.

\propoptions{counter=P, style=proposition}

\begin{prop}
  \item First~P. \label{po:P1}
  \item Second~P. \label{po:P2}
\end{prop}

\begin{prop}
  \item Third~P (continues from previous). \label{po:P3}
  \item Fourth~P. \label{po:P4}
\end{prop}

\propoptions{counter=numpropi, style=numbered}   % restore

Refs: \ref{po:P1}, \ref{po:P2}, \ref{po:P3}, \ref{po:P4}.

\subsection{Env-arg overrides propoptions}

% \propoptions{counter=P} sets P as default.
% The second env overrides with counter=numpropi for that env only.

\propoptions{counter=P, style=proposition}

\begin{prop}
  \item P-numbered (from propoptions). \label{po:ov1}
\end{prop}
\begin{prop}[counter=numpropi, style=numbered]
  \item numpropi-numbered (env overrides propoptions). \label{po:ov2}
\end{prop}
\begin{prop}
  \item Back to P (propoptions restored). \label{po:ov3}
\end{prop}

\propoptions{counter=numpropi, style=numbered}  % restore

Refs: \ref{po:ov1} (P), \ref{po:ov2} (numpropi), \ref{po:ov3} (P).

\subsection{Group-scoped \texttt{\textbackslash propoptions}}

% Inside the group, counter=P; outside, the default numpropi is restored.

Before group: numpropi style.
\begin{prop}
  \item Numbered item (numpropi). \label{po:gr0}
\end{prop}

\begingroup
  \propoptions{counter=P, style=proposition}
  Inside group: P style.
  \begin{prop}
    \item P-numbered (group scope). \label{po:gr1}
    \item P-numbered (group scope). \label{po:gr2}
  \end{prop}
\endgroup

After group: numpropi restored.
\begin{prop}
  \item Numbered item (numpropi again). \label{po:gr3}
\end{prop}

Refs: \ref{po:gr0}, \ref{po:gr1}, \ref{po:gr2}, \ref{po:gr3}.

% ============================================================
\section{Global default style changes}
% ============================================================

% named style= changes which style is used for items with a given name
% (but no explicit style=).

\propoptions{named style=thesis}
\begin{prop}
  \item[First Claim] Should now be thesis style (small-caps, flush). \label{gs:A}
  \item[Second Claim] Also thesis. \label{gs:B}
\end{prop}
\propoptions{named style=proposition}   % restore default

Refs: \ref{gs:A}, \ref{gs:B}.

\end{document}
